% ---------------------------------------------------------------- Приложение А
\appendixsection{А}{справочное}{Справочник функций Python для лабораторных работ}
\label{app:python}

\begin{table}[!htb]
\caption{Функции NumPy, SciPy и библиотек машинного обучения, используемые в практикуме}
\label{tab:python-ref}
\centering\small
\begin{tabularx}{\textwidth}{>{\raggedright}p{4.0cm}>{\raggedright\arraybackslash}X}
\toprule
Задача & Функции \\
\midrule
Случайные числа (воспроизводимо) & \code{rng = np.random.default\_rng(seed)}; \code{rng.random}, \code{rng.integers}, \code{rng.normal}, \code{rng.uniform}, \code{rng.choice(p=...)}, \code{rng.permutation} \\
Векторные операции & \code{@}, \code{np.outer}, \code{np.where}, \code{np.argsort}, \code{np.argmin}, \code{np.argmax}, \code{np.roll}, \code{np.cumsum}, \code{np.searchsorted}, \code{np.bitwise\_xor.accumulate} \\
Линейная алгебра & \code{np.linalg.lstsq}, \code{np.linalg.pinv}, \code{np.linalg.solve}, \code{np.linalg.eigvalsh} \\
Временные ряды & \code{np.lib.stride\_tricks.sliding\_window\_view} \\
Предобработка & \code{sklearn.preprocessing.MinMaxScaler}, \code{sklearn.model\_selection.train\_test\_split} \\
Модели для сравнения & \code{Perceptron}, \code{LinearRegression}, \code{SGDRegressor}, \code{MLPRegressor} (\code{sklearn}); \code{torch.nn.Sequential}, \code{torch.optim}; \code{MiniSom}; \code{KMeans} \\
Метрики & \code{mean\_squared\_error}, \code{r2\_score}, \code{adjusted\_rand\_score} (\code{sklearn.metrics}); \code{som.quantization\_error}, \code{som.topographic\_error} \\
Эталонные решения и проверка & \code{scipy.optimize.minimize\_scalar}, \code{differential\_evolution}, \code{least\_squares(method='lm')}, \code{check\_grad} \\
Эволюционные библиотеки & \code{deap.tools} (\code{selTournament}, \code{cxUniform}, \code{cxBlend}, \code{cxPartialyMatched}, \code{cxOrdered}); \code{pygad.GA} \\
Коммивояжёр & \code{scipy.spatial.distance.cdist}, \code{np.triu\_indices}, пакет \code{tsplib95} \\
Безградиентная оптимизация & \code{cma.CMAEvolutionStrategy}, \code{cma.fmin2}; \code{optuna.create\_study}, \code{TPESampler}, \code{CmaEsSampler}, \code{NSGAIISampler}; \code{pymoo}: \code{NSGA2}, \code{minimize}, \code{get\_problem("zdt1")}, \code{HV}, \code{IGD} \\
Подбор гиперпараметров & \code{sklearn.model\_selection.cross\_val\_score}, \code{StratifiedKFold}, \code{GridSearchCV}, \code{RandomizedSearchCV}; \code{sklearn.datasets.load\_wine}, \code{load\_digits}, \code{load\_breast\_cancer} \\
Проверка и запуск кода модели & \code{ast.parse}, \code{ast.walk}, \code{subprocess.run(..., timeout=...)}, \code{concurrent.futures.ThreadPoolExecutor}; \code{anthropic.Anthropic().messages.create} \\
Графики и анимация & \code{plt.contourf}, \code{plt.scatter}, \code{matplotlib.animation.FuncAnimation}, \code{plotly.graph\_objects.FigureWidget} \\
Интерфейс в Jupyter & \code{ipywidgets}: \code{IntSlider}, \code{FloatText}, \code{Dropdown}, \code{Button}, \code{Play}, \code{interact} \\
\bottomrule
\end{tabularx}
\end{table}

% ---------------------------------------------------------------- Приложение Б
\appendixsection{Б}{справочное}{Листинги программ}
\label{app:listings}

\emph{Операторы для задачи коммивояжёра (ЛР~№7)}: чтение файла TSPLIB, матрица расстояний
\code{EUC\_2D}, циклический кроссинговер и векторный локальный поиск 2-opt; последняя строка проверяет
CX на примере подраздела~\ref{sec:tsp}.

\lstinputlisting{listings/tsp_cx.py}

\emph{CMA-ES и IPOP-CMA-ES (ЛР~№8)} по учебнику Н.~Хансена~\cite{hansen2016}: порождение точек,
взвешенная рекомбинация, пути эволюции, ранг-1 и ранг-$\mu$ обновления ковариационной матрицы,
кумулятивная адаптация шага.

\lstinputlisting{listings/cma_es.py}

\emph{Мини-FunSearch (ЛР~№9)}: генерация примеров и оценка эвристики упаковки, классические правила,
извлечение и статическая проверка кода от языковой модели, оценка в отдельном процессе с тайм-аутом,
вызов модели через Anthropic SDK или Claude Code CLI и эволюционный цикл.

\lstinputlisting{listings/funsearch.py}
